\scanpage{55}
On dit alors que le morphisme \deleted{\unsure{fini}} plat
$T\unsure{\mathbin{-}|\to} S$ est \uline{homogène} \inserted{si $\unsure{\nu|v}(s)$ est loc. constant}. Cette notion est [dans le
cas de morphismes \uline{finis} du moins]
invariante par composition des morphismes,
extension de la base [trivial sur (i), qui
sera la définition]

\deleted{Corollaire}

\deleted{\uline{Proposition} Soient \deleted{$\unsure{X,Y,S}$} $T,S,X,Y$ comme
dans \unsure{ds} de Prop. précédente. Supposons
que $X$ soit un schéma en groupes sur $Y$,
et $P$ un fibré principal sur $X$}

\uline{Cor. 4} Sous les conditions de (i), supposons
$T/S$ fini, \deleted{et que $f$ \unsure{soit} \illegible{1} \unsure{l'élément}} \inserted{Si $g$ est surjectif,}
Alors $f_\ast(g)$ est \uline{surjectif}

\uline{Proposition} Le foncteur $f_\ast$ transforme
fibrés \inserted{\unsure{formellement|formellt|formelles}} principaux sur un $Y/T$
[sous réserve d'existence \inserted{et sans autre restriction}] en
fibrés \unsure{formellement|formellt|formelles} principaux.
% slanted
\margin{Pour éviter mention \unsure{compliquée} \unsure{trivialité des solutions|triviale des solutions|principal homogène}}
Il transforme fibrés principaux
simples en fibrés principaux
simples, sans supposer que $T/S$ soit fini.
Il transforme \deleted{\unsure{fibrés}} fibrés principaux à groupes
\unsure{et} \deleted{\unsure{finis|fermés}} simples \unsure{pour séparables|sont représentables|séparables} en \deleted{\unsure{des}} des
% archive sheet number: 54
